\subsection{Metodologia propuesta}

Antes de la carga tensorial, las imágenes del repositorio se someterán a una verificación estricta de formato e integridad mediante la biblioteca PIL. El algoritmo filtra de manera automática cualquier archivo corrupto o con extensiones distintas a .jpg, .jpeg y .png. Posteriormente, los tensores de entrada se redimensionan al tamaño estándar de la arquitectura de destino mediante torchvision. Finalmente, las matrices de píxeles se convierten al rango flotante de 0 a 1 y se estandarizan por canal utilizando la media ($0.485, 0.456, 0.406$) y desviación estándar ($0.229, 0.224, 0.225$) del conjunto ImageNet.

El particionamiento del conjunto de datos en las carpetas train, val y test se respeta de manera prioritaria y estricta desde el sistema de archivos. Esta separación física se ejecuta de forma previa a cualquier operación de transformación o carga en memoria. Los subconjuntos de validación y prueba emplean únicamente una transformación de evaluación determinista sin distorsiones estocásticas. Con este diseño metodológico riguroso se elimina completamente la posibilidad de fuga de datos durante el entrenamiento.

El proceso de expansión sintética de datos se aplica exclusivamente al conjunto de entrenamiento durante el flujo de trabajo optimizado. Las transformaciones de aumento en PyTorch incluyen giros horizontales aleatorios y rotaciones espaciales en un rango de $\pm 10^{\circ}$. Asimismo, se incorporan perturbaciones fotométricas dinámicas en el espacio cromático mediante ColorJitter con variaciones de $0.1$ en brillo y contraste. Estas modificaciones en tiempo de ejecución incrementan la invariancia de la red ante artefactos visuales sin alterar la pureza del conjunto de test.

El pipeline computacional se ejecuta utilizando la función de pérdida CrossEntropyLoss y el optimizador Adam sobre la GPU de Google Colab. En la Fase 2 optimizada, se implementa precisión mixta automática (AMP con GradScaler) para acelerar el tiempo de cómputo. El ajuste de hiperparámetros se controla dinámicamente mediante el programador ReduceLROnPlateau y decaimiento de pesos. Finalmente, se activa un stopper de convergencia basado en Early Stopping sobre el val\_loss para preservar automáticamente los pesos del mejor modelo.